% EXAM PRACTICE preamble: copied from ../worksheet/ws_preamble.tex (2026-10-10).
% Changes: header text; \Q takes a section tag; student copy has NO working
% boxes (questions follow one another); teacher copy: one question per page,
% solution box below the question.
% Original note: Worksheet preamble: LOCAL COPIES (ws_preamble_build, ws_preamble_local,
% ws_figstyles; copied from ../notes on 2026-10-09) + chapter styles,
% plus RI-style question / solution-box layout (item 16, item 79).
% WORKSHEET COPY of ../notes/preamble_build.tex (2026-10-09 version), unchanged.
% preamble_build.tex: copy of Maximus's preamble.tex for the DE v3 build.
% The blocks marked [BUILD COPY CHANGE] are the earlier build fixes.
% The blocks marked [V3 ADDITION] are new in this version (example box,
% spacing). Everything else is unchanged from preamble.tex.

% For margins
\usepackage[a4paper, margin = 20mm]{geometry}

\usepackage[Lenny]{fncychap}
% Rejne, Glenn, Lenny

\usepackage{etoc}

% [V3 ADDITION] cancel: used to show the dx's "cancelling" in Section 11.3
\usepackage{cancel}
\renewcommand{\CancelColor}{\color{red}}

% For enumeration & math symbols
\usepackage{enumerate}
\usepackage{amsmath,amsfonts,amssymb,amsthm,graphicx,lipsum}

% For diagrams, plots and boxes
\usepackage{pgfplots}
\pgfplotsset{compat=1.18}
\usepackage{tikz}
\usepackage[most]{tcolorbox}
\tcbuselibrary{skins}
\usepackage{xcolor}
\pgfplotsset{
    standard/.style={clip = false,
    axis line style = thick,
    trig format=rad,
    axis lines = middle,
    every axis x label/.style={at={(current axis.right of origin)},anchor= west},
    every axis y label/.style={at={(current axis.above origin)},anchor=south}
    }
}

% For GC font
\usepackage{fontspec}
% [BUILD COPY CHANGE 1] The TI-84 key font is loaded only if the file exists
% (checked in ./Fonts/... and ../Fonts/...). Otherwise \ti{} falls back to a
% small boxed sans-serif key. Known letters of the TI font are mapped to their
% key names in the fallback: o = [y=], z = [mode], f = [L3].
\newcommand{\tikey}[1]{\tcbox[on line, enhanced, boxsep=0pt, left=2pt, right=2pt,
  top=1pt, bottom=1pt, boxrule=0.5pt, arc=1.5pt, colback=white,
  colframe=black!70]{\sffamily\footnotesize\strut #1}}
\ExplSyntaxOn
\NewDocumentCommand{\tifallback}{m}{
  \str_case:nnF {#1}
    { {o}{\tikey{y=}} {z}{\tikey{mode}} {f}{\tikey{L3}} }
    { \tikey{#1} }
}
\ExplSyntaxOff
\def\tifontcase{0}
\IfFileExists{Fonts/TI-84 Plus CE Key Fonts/TI84PlusCEKeys.ttf}
  {\def\tifontcase{1}}{}
\IfFileExists{../Fonts/TI-84 Plus CE Key Fonts/TI84PlusCEKeys.ttf}
  {\def\tifontcase{2}}{}
\ifcase\tifontcase\relax
  % 0: font not found, use boxed fallback keys
  \newcommand{\ti}[1]{\tifallback{#1}}
\or
  % 1: original location (as in preamble.tex)
  \newfontfamily\tifont[
    Path = {Fonts/TI-84 Plus CE Key Fonts/},
    UprightFont = {TI84PlusCEKeys.ttf}
  ]{}
  \newcommand{\ti}[1]{{\tifont #1}}
\or
  % 2: one level up (when compiling inside draft-ch1/)
  \newfontfamily\tifont[
    Path = {../Fonts/TI-84 Plus CE Key Fonts/},
    UprightFont = {TI84PlusCEKeys.ttf}
  ]{}
  \newcommand{\ti}[1]{{\tifont #1}}
\fi

% [AoI items 17/18/37] GC DISPLAY FONT. Every GC key, command, menu item and
% screen line goes through \gcfont, so the font can be swapped HERE in one place.
%   \gc{fnInt(}      command / menu item / stored variable in the display font
%   \gckey{2nd}      a calculator key (display font in a small key outline)
%   \gcscreen{...}   a calculator screen (display font on a grey-green LCD panel)
% The font file lives in ../Fonts (shared by notes, slides and worksheet).
\def\gcfontfile{VT323-Regular.ttf}
\IfFileExists{../Fonts/\gcfontfile}{%
  \newfontfamily\gcfont[Path=../Fonts/, Scale=1.32]{\gcfontfile}}{%
  \newcommand\gcfont{\ttfamily}}
\newcommand{\gc}[1]{\mbox{\gcfont #1}}
\newcommand{\gckey}[1]{\tcbox[on line, enhanced, boxsep=0pt, left=2pt, right=2pt,
  top=0.6pt, bottom=0.2pt, boxrule=0.5pt, arc=1.5pt, colback=black!4,
  colframe=black!65]{\gcfont\strut #1}}
\definecolor{lcd}{RGB}{214,222,204}
\newcommand{\gcscreen}[1]{\tcbox[on line, enhanced, boxsep=1pt, left=4pt, right=4pt,
  top=2pt, bottom=2pt, boxrule=1.2pt, arc=1pt, colback=lcd, colframe=black!70]{%
  \gcfont\begin{tabular}{@{}l@{}}#1\end{tabular}}}
% integral template as shown on the TI-84 Plus CE (MathPrint): \gcint{0}{2}{e^(-X²)}{X}
\newcommand{\gcint}[4]{{\gcfont\raisebox{-0.35ex}{\Large ∫}\hspace{-0.15em}%
  \raisebox{-0.2ex}{\begin{tabular}{@{}l@{}}\footnotesize #2\\[-1.3ex]\footnotesize #1\end{tabular}}%
  \hspace{0.1em}(#3)d#4}}
% old key macro kept as an alias so nothing breaks
\renewcommand{\ti}[1]{\gckey{#1}}

% SI units
\usepackage{siunitx}

% setspace
%\usepackage{setspace}
%\onehalfspacing

% For hyperlinks within and out of document
\usepackage{hyperref}
\hypersetup{
    colorlinks = true,
    linkcolor = black,
    urlcolor = blue,
}   

% For captions for pics
\usepackage{capt-of}


% tcolorbox custom environments
% Definition ========================================================
\newtcolorbox[auto counter,number within=section]{definition}[2][drop shadow]{%
enhanced, before skip=5mm, after skip=6mm, % [V3 ADDITION] spacing
colback=blue!5!white,colframe=blue!50!black,fonttitle=\bfseries,
title= Definition~\thetcbcounter. #2,#1}

% Theorem ==========================================================
\newtcolorbox[auto counter,number within=section]{theorem}[2][drop shadow]{%
enhanced, before skip=5mm, after skip=6mm, % [V3 ADDITION] spacing
colback=black!5!white,colframe=black,fonttitle=\bfseries,
title= Theorem~\thetcbcounter. #2,#1}

% Caution box ======================================================
\newtcolorbox{caution}[1][]{enhanced,
before skip=5mm,after skip=6mm, % [V3 ADDITION] was 2mm/3mm
boxrule=1pt,left=11mm,right=2mm,top=3mm,bottom=3mm,
colback=yellow!50,
colframe=yellow!20!black,
sharp corners,rounded corners=southeast,arc is angular,arc=3mm,
underlay={%
\path[fill=tcbcolback!80!black] ([yshift=3mm]interior.south east)--++(-0.4,-0.1)--++(0.1,-0.2);
\path[draw=tcbcolframe,shorten <=-0.05mm,shorten >=-0.05mm] ([yshift=3mm]interior.south east)--++(-0.4,-0.1)--++(0.1,-0.2);
\path[fill=yellow!50!black,draw=none] (interior.south west) rectangle node[white]{\Huge\bfseries !} ([xshift=9mm]interior.north west);
},
drop fuzzy shadow,#1}

% Question box ======================================================
\newtcolorbox{question}[1][]{enhanced,
before skip=5mm,after skip=6mm, % [V3 ADDITION] was 2mm/3mm
boxrule=1pt,left=11mm,right=2mm,top=3mm,bottom=3mm,
colback=blue!20!white,
colframe=blue!20!black,
sharp corners,rounded corners=southeast,arc is angular,arc=3mm,
underlay={%
\path[fill=tcbcolback!80!black] ([yshift=3mm]interior.south east)--++(-0.4,-0.1)--++(0.1,-0.2);
\path[draw=tcbcolframe,shorten <=-0.05mm,shorten >=-0.05mm] ([yshift=3mm]interior.south east)--++(-0.4,-0.1)--++(0.1,-0.2);
\path[fill=blue!50!black,draw=none] (interior.south west) rectangle node[white]{\Huge\bfseries ?} ([xshift=9mm]interior.north west);
},
drop fuzzy shadow,#1}

% Information box ======================================================
\newtcolorbox{info}[1][]{enhanced,
before skip=5mm,after skip=6mm, % [V3 ADDITION] was 2mm/3mm
boxrule=1pt,left=11mm,right=2mm,top=3mm,bottom=3mm,
colback=cyan!20!white,
colframe=cyan!20!black,
sharp corners,rounded corners=southeast,arc is angular,arc=3mm,
underlay={%
\path[fill=tcbcolback!80!black] ([yshift=3mm]interior.south east)--++(-0.4,-0.1)--++(0.1,-0.2);
\path[draw=tcbcolframe,shorten <=-0.05mm,shorten >=-0.05mm] ([yshift=3mm]interior.south east)--++(-0.4,-0.1)--++(0.1,-0.2);
\path[fill=cyan!50!black,draw=none] (interior.south west) rectangle node[white]{\Huge\bfseries i} ([xshift=9mm]interior.north west);
},
drop fuzzy shadow,#1}

% Mistake box ======================================================
\newtcolorbox{mistake}[1][]{enhanced,
before skip=5mm,after skip=6mm, % [V3 ADDITION] was 2mm/3mm
boxrule=1pt,left=11mm,right=2mm,top=3mm,bottom=3mm,
colback=red!20!white,
colframe=red!20!black,
sharp corners,rounded corners=southeast,arc is angular,arc=3mm,
underlay={%
\path[fill=tcbcolback!80!black] ([yshift=3mm]interior.south east)--++(-0.4,-0.1)--++(0.1,-0.2);
\path[draw=tcbcolframe,shorten <=-0.05mm,shorten >=-0.05mm] ([yshift=3mm]interior.south east)--++(-0.4,-0.1)--++(0.1,-0.2);
\path[fill=red!50!black,draw=none] (interior.south west) rectangle node[white]{\Huge\bfseries $\times$} ([xshift=9mm]interior.north west);
},
drop fuzzy shadow,#1}



% [V3 ADDITION] Worked example box ==================================
% Own counter, numbered within the section (Example 11.2.1, 11.2.2, ...).
% For per-subsection numbering (11.2.1.1) change `section' to `subsection'. Blue `question' boxes are kept for thought-provoking questions
% only; worked examples always use this box. Usage:
%   \begin{example} ...question text... \end{example}
%   \solution{...}
\newtcolorbox[auto counter,number within=section]{example}[1][]{%
enhanced, breakable=false, before skip=5mm, after skip=4mm,
colback=teal!5!white, colframe=teal!55!black, boxrule=0.8pt,
left=3mm, right=3mm, top=2mm, bottom=2mm,
sharp corners, rounded corners=northwest, arc=3mm,
fonttitle=\bfseries, coltitle=white,
title=Example~\thetcbcounter, attach boxed title to top left={yshift=-2mm, xshift=3mm},
boxed title style={colback=teal!55!black, sharp corners, rounded corners=northwest, arc=2mm},
drop fuzzy shadow, #1}

% [V3 ADDITION] Spacing between sections, subsections, parts, paragraphs =
% (fncychap's chapter heading is untouched; only \section and below.)
\makeatletter
\renewcommand\section{\@startsection{section}{1}{\z@}%
  {-2ex \@plus -.5ex \@minus -.2ex}{3.2ex \@plus .3ex}%
  {\normalfont\LARGE\bfseries\color{blue!50!black}}}
\renewcommand\subsection{\@startsection{subsection}{2}{\z@}%
  {-1.5ex \@plus -.5ex \@minus -.2ex}{2.6ex \@plus .3ex}%
  {\normalfont\Large\bfseries}}
\renewcommand\subsubsection{\@startsection{subsubsection}{3}{\z@}%
  {-4ex \@plus -1ex \@minus -.5ex}{1.6ex \@plus .2ex}%
  {\normalfont\large\bfseries\color{teal!50!black}}}
\makeatother
\setlength{\parindent}{0pt}
\setlength{\parskip}{0.9ex plus 0.3ex minus 0.2ex}

% Headers and Footers
\usepackage{fancyhdr,lastpage}
% [V3 R2] grey left-margin line on the out-of-syllabus pages (item 13)
\usepackage{eso-pic}
% [AoI item 63] thicker line (2.6pt, was 1.6pt) that starts just BELOW the header rule
% (top 14mm of the page left clear). Flag: copy this line into every chapter's preamble.
\newcommand{\oosmargin}{\AddToShipoutPictureBG{\put(12mm,0){\color{gray!75}\rule{2.6pt}{\dimexpr\paperheight-14mm\relax}}}}
\pagestyle{fancy}
\lhead{Maximus Pung Jun}
% [BUILD COPY CHANGE 2] Header bug fix: \chaptername is the word "Chapter",
% so the original header printed "Chapter 1: Chapter". \chaptermark now stores
% the bare chapter title in \leftmark, and the header uses \leftmark.
\renewcommand{\chaptermark}[1]{\markboth{#1}{}}
\rhead{Chapter \thechapter: \leftmark}
% \cfoot{Page \thepage\  of \pageref{LastPage}} % buggy, does not appear on chapter pages.

% [V3 R2] tick and cross markers for mistake boxes (G3)
\newcommand{\wrong}{\textcolor{red!70!black}{\ensuremath{\boldsymbol{\times}}}}
\newcommand{\cmark}{\textcolor{green!50!black}{\ensuremath{\checkmark}}}

\newcommand{\xmark}{\wrong}   % alias: red cross
% [V3 R2] out-of-syllabus helpers (DE notes item 13)
%   \oosmargin    : grey line in the left margin on every following page
%   \oosmarginoff : stop the margin line
%   \oosnote      : grey banner "Out of syllabus. For your interest only."
%   outofsyllabus : environment = new page + margin line + banner; ends with
%                   \clearpage and switches the margin line off again
\newcommand{\oosmarginoff}{\ClearShipoutPictureBG}
\newcommand{\oosnote}[1][Nothing in this section is examined in H2 Mathematics.]{\par\noindent\colorbox{gray!15}{\parbox{\dimexpr\textwidth-2\fboxsep}{\textbf{Out of syllabus. For your interest only.} #1}}\par}
\newenvironment{outofsyllabus}{\clearpage\oosmargin}{\clearpage\oosmarginoff}

% Notation as in Maximus's notes: italic d, display style everywhere (G1).
% Fractions inside exponents stay small, so never use these macros inside ^{}.
\newcommand{\dydx}{\dfrac{dy}{dx}}
\newcommand{\dvdx}{\dfrac{dv}{dx}}
\newcommand{\dudx}{\dfrac{du}{dx}}
\newcommand{\ddydx}{\dfrac{d^{2}y}{dx^{2}}}
\newcommand{\ddt}[1]{\dfrac{d#1}{dt}}
% Small syllabus flag (no longer used in the DE notes, kept for other chapters)
\newcommand{\syllabusflag}[1]{\par\noindent{\footnotesize\color{red!70!black}\textbf{Syllabus note.} #1}\par}


% Teacher's copy toggle ================================================
% [BUILD COPY CHANGE 3] \ifx compares the whole word (the original
% \if\conform test was always false). DE_notes_student.tex predefines
% \conform as HideSolutions.
% A student copy defines \conform as HideSolutions BEFORE loading this preamble.
\providecommand\conform{ShowSolutions}
\newcommand\ShowSolutionsKey{ShowSolutions}
% [V3 ADDITION] In the student copy the hidden solution leaves stretchable
% blank space (\vfill), so the free space on each page is shared out as
% working space below the examples instead of collecting at the bottom.
% [AoI item 10] "Solution." on its own line, so that (a), (b), (c) line up below it.
\newcommand\solution[1]{\ifx\conform\ShowSolutionsKey
  \par\noindent\textbf{Solution.}\par\nobreak\vspace{0.5ex}\noindent #1\par\medskip\else\par\vspace{\workspace}\fi}
% [V3 R2] minimum working space left under each example in the student copy
\newlength\workspace\setlength\workspace{3cm plus 1fill}


\newcommand{\NN}{\mathbb{N}}
\newcommand{\ZZ}{\mathbb{Z}}
\newcommand{\QQ}{\mathbb{Q}}
\newcommand{\RR}{\mathbb{R}}
\newcommand{\CC}{\mathbb{C}}


% WORKSHEET COPY of ../notes/preamble_local.tex (version the worksheet was last built with,
% from ../_backup_pre_limitsum, 2026-10-09 17:31). Local so the worksheet no longer depends
% on ../notes; changes: % WORKSHEET COPY of ../notes/figstyles.tex (2026-10-09 17:31 version). Only change (follow-up):
% thin light axes in `sketch' and `ax3' (marked [2026-10-09 follow-up]).
% =====================================================================
% figstyles.tex: plot styles and figure macros shared by the notes, the
% slides (../slides) and the worksheets (../worksheet). Needs pgfplots,
% the fillbetween library and tikz libraries arrows.meta and calc.
% House rules built in:
%   * `axis on top' in every sketch: shading never covers an axis, and the
%     fill starts at the axis line (items 31, 32, 34);
%   * regions are filled from the curve itself (fill between / closed plot);
%   * 3D solids in the page-13 style (translucent blue, light end caps).
% =====================================================================
\pgfplotsset{
  sketch/.style={clip=false, axis line style={thin, black!45}, every axis label/.append style={black!60}, trig format=rad, % [2026-10-09 follow-up] was axis line style=thick
    axis lines=middle, axis on top,
    every axis x label/.style={at={(current axis.right of origin)},anchor=west},
    every axis y label/.style={at={(current axis.above origin)},anchor=south},
    ticks=none, xlabel=$x$, ylabel=$y$, width=7.6cm, height=6cm,
    every node/.append style={font=\footnotesize}},
  curve/.style={blue, semithick, smooth, samples=150},
  line2/.style={red!75!black, semithick, samples=2},
  % 3D solids of revolution (angles in DEGREES: no `trig format=rad' here)
  solid3d/.style={hide axis, trig format=deg, width=8.4cm, height=6.4cm, clip=false,
    z buffer=sort, colormap={bl}{rgb255=(150,185,235) rgb255=(40,90,180)},
    every node/.append style={font=\footnotesize}},
  solidsurf/.style={surf, shader=interp, opacity=0.75, point meta=y},
  red3d/.style={surf, shader=flat, draw=none, point meta=1,
    colormap={r}{rgb255=(215,50,50) rgb255=(215,50,50)}},
  orange3d/.style={surf, shader=interp, point meta=z,
    colormap={o}{rgb255=(250,210,160) rgb255=(220,120,30)}},
  cap3d/.style={surf, shader=flat, draw=none, point meta=1,
    colormap={c}{rgb255=(185,205,238) rgb255=(185,205,238)}},
  mini3d/.style={solid3d, width=2.6cm, height=2.4cm},
}
\tikzset{
  region/.style={fill=blue!22, draw=none},
  regionB/.style={fill=orange!35, draw=none},
  strip/.style={fill=red!55, draw=red!70!black, thin},
  asym/.style={black, dashed, thin},
  pt/.style={circle, fill=black, inner sep=1.1pt},
  dimarrow/.style={<->, >=Stealth, thin},
  ax3/.style={-{Stealth[length=1.5mm]}, thin, black!50}, % [2026-10-09 follow-up] was -{Stealth[length=2mm]}, thick
}
\definecolor{hl}{RGB}{30,110,200}

% ---------------------------------------------------------------------
% Mini 3D solids for SOLUTIONS (items 54, 62, 66, 74): solid only, no
% washer, every picture 1.5 cm tall; a "big minus small" pair is 4.4 cm wide.
%  \solX[x unit cm]{R(x)}{a}{b}{R(a)}{R(b)}{x-axis end}{y-axis top}{curve colour}
%     solid swept out by y = R(x), a <= x <= b, about the x-axis
%  \solY[y unit cm]{R(y)}{c}{d}{R(c)}{R(d)}{x-axis end}{y-axis top}{curve colour}
%     solid swept out by x = R(y), c <= y <= d, about the y-axis
%     (R is written in terms of the letter x, which here stands for y)
% ---------------------------------------------------------------------
\newcommand{\solX}[9][0.5]{%
\begin{tikzpicture}[baseline=(current bounding box.center)]
\begin{axis}[mini3d, x={(#1cm,-0.06cm)}, y={(0.12cm,0.09cm)}, z={(0cm,0.4cm)}]
  \addplot3[cap3d, domain=0:1, y domain=0:360, samples=2, samples y=40] ({#3}, {x*(#5)*cos(y)}, {x*(#5)*sin(y)});
  \addplot3[solidsurf, opacity=0.5, point meta=z, domain=#3:#4, y domain=0:360, samples=24, samples y=40] ({x}, {(#2)*cos(y)}, {(#2)*sin(y)});
  \addplot3[cap3d, domain=0:1, y domain=0:360, samples=2, samples y=40] ({#4}, {x*(#6)*cos(y)}, {x*(#6)*sin(y)});
  \addplot3[#9, thick, domain=#3:#4, samples=40, samples y=1] ({x}, {0}, {#2});
  \draw[ax3] (axis cs:{min(0,#3)-0.3},0,0) -- (axis cs:#7,0,0);
  \draw[ax3] (axis cs:0,0,{-0.15*(#8)}) -- (axis cs:0,0,#8);
\end{axis}
\end{tikzpicture}}
\newcommand{\solY}[9][0.6]{%
\begin{tikzpicture}[baseline=(current bounding box.center)]
\begin{axis}[mini3d, x={(0.3cm,-0.03cm)}, y={(0.13cm,0.11cm)}, z={(0cm,#1cm)}]
  \addplot3[cap3d, domain=0:1, y domain=0:360, samples=2, samples y=40] ({x*(#5)*cos(y)}, {x*(#5)*sin(y)}, {#3});
  \addplot3[solidsurf, opacity=0.5, point meta=z, domain=#3:#4, y domain=0:360, samples=24, samples y=40] ({(#2)*cos(y)}, {(#2)*sin(y)}, {x});
  \addplot3[cap3d, domain=0:1, y domain=0:360, samples=2, samples y=40] ({x*(#6)*cos(y)}, {x*(#6)*sin(y)}, {#4});
  \addplot3[#9, thick, domain=#3:#4, samples=40, samples y=1] ({#2}, {0}, {x});
  \draw[ax3] (axis cs:0,0,{min(0,#3)-0.15*(#8)}) -- (axis cs:0,0,#8);
  \draw[ax3] (axis cs:-0.3,0,0) -- (axis cs:#7,0,0);
\end{axis}
\end{tikzpicture}}
\newcommand{\minisolo}[1]{\adjustbox{totalheight=1.5cm, max width=4.6cm}{#1}}
\newcommand{\minipair}[2]{\adjustbox{totalheight=1.5cm, max width=4.6cm}{#1\;\raisebox{0.15ex}{\Large$-$}\;#2}}
\newcommand{\minilabel}[1]{\raisebox{0.65cm}{\small #1}\;}
; [2026-10-09 follow-up] \mk marks now bold, as in the notes.
% =====================================================================
% preamble_local.tex: extra packages, styles and macros for the
% Definite Integrals chapter ONLY. Loaded after preamble_build.tex.
% Plot styles and the mini 3D macros are in figstyles.tex (shared with the
% slides and worksheets).
% =====================================================================
\usepgfplotslibrary{fillbetween}
\usepackage{adjustbox} % mini 3D pictures (figstyles.tex)
\usetikzlibrary{arrows.meta, decorations.pathreplacing, calc}
% WORKSHEET COPY of ../notes/figstyles.tex (2026-10-09 17:31 version). Only change (follow-up):
% thin light axes in `sketch' and `ax3' (marked [2026-10-09 follow-up]).
% =====================================================================
% figstyles.tex: plot styles and figure macros shared by the notes, the
% slides (../slides) and the worksheets (../worksheet). Needs pgfplots,
% the fillbetween library and tikz libraries arrows.meta and calc.
% House rules built in:
%   * `axis on top' in every sketch: shading never covers an axis, and the
%     fill starts at the axis line (items 31, 32, 34);
%   * regions are filled from the curve itself (fill between / closed plot);
%   * 3D solids in the page-13 style (translucent blue, light end caps).
% =====================================================================
\pgfplotsset{
  sketch/.style={clip=false, axis line style={thin, black!45}, every axis label/.append style={black!60}, trig format=rad, % [2026-10-09 follow-up] was axis line style=thick
    axis lines=middle, axis on top,
    every axis x label/.style={at={(current axis.right of origin)},anchor=west},
    every axis y label/.style={at={(current axis.above origin)},anchor=south},
    ticks=none, xlabel=$x$, ylabel=$y$, width=7.6cm, height=6cm,
    every node/.append style={font=\footnotesize}},
  curve/.style={blue, semithick, smooth, samples=150},
  line2/.style={red!75!black, semithick, samples=2},
  % 3D solids of revolution (angles in DEGREES: no `trig format=rad' here)
  solid3d/.style={hide axis, trig format=deg, width=8.4cm, height=6.4cm, clip=false,
    z buffer=sort, colormap={bl}{rgb255=(150,185,235) rgb255=(40,90,180)},
    every node/.append style={font=\footnotesize}},
  solidsurf/.style={surf, shader=interp, opacity=0.75, point meta=y},
  red3d/.style={surf, shader=flat, draw=none, point meta=1,
    colormap={r}{rgb255=(215,50,50) rgb255=(215,50,50)}},
  orange3d/.style={surf, shader=interp, point meta=z,
    colormap={o}{rgb255=(250,210,160) rgb255=(220,120,30)}},
  cap3d/.style={surf, shader=flat, draw=none, point meta=1,
    colormap={c}{rgb255=(185,205,238) rgb255=(185,205,238)}},
  mini3d/.style={solid3d, width=2.6cm, height=2.4cm},
}
\tikzset{
  region/.style={fill=blue!22, draw=none},
  regionB/.style={fill=orange!35, draw=none},
  strip/.style={fill=red!55, draw=red!70!black, thin},
  asym/.style={black, dashed, thin},
  pt/.style={circle, fill=black, inner sep=1.1pt},
  dimarrow/.style={<->, >=Stealth, thin},
  ax3/.style={-{Stealth[length=1.5mm]}, thin, black!50}, % [2026-10-09 follow-up] was -{Stealth[length=2mm]}, thick
}
\definecolor{hl}{RGB}{30,110,200}

% ---------------------------------------------------------------------
% Mini 3D solids for SOLUTIONS (items 54, 62, 66, 74): solid only, no
% washer, every picture 1.5 cm tall; a "big minus small" pair is 4.4 cm wide.
%  \solX[x unit cm]{R(x)}{a}{b}{R(a)}{R(b)}{x-axis end}{y-axis top}{curve colour}
%     solid swept out by y = R(x), a <= x <= b, about the x-axis
%  \solY[y unit cm]{R(y)}{c}{d}{R(c)}{R(d)}{x-axis end}{y-axis top}{curve colour}
%     solid swept out by x = R(y), c <= y <= d, about the y-axis
%     (R is written in terms of the letter x, which here stands for y)
% ---------------------------------------------------------------------
\newcommand{\solX}[9][0.5]{%
\begin{tikzpicture}[baseline=(current bounding box.center)]
\begin{axis}[mini3d, x={(#1cm,-0.06cm)}, y={(0.12cm,0.09cm)}, z={(0cm,0.4cm)}]
  \addplot3[cap3d, domain=0:1, y domain=0:360, samples=2, samples y=40] ({#3}, {x*(#5)*cos(y)}, {x*(#5)*sin(y)});
  \addplot3[solidsurf, opacity=0.5, point meta=z, domain=#3:#4, y domain=0:360, samples=24, samples y=40] ({x}, {(#2)*cos(y)}, {(#2)*sin(y)});
  \addplot3[cap3d, domain=0:1, y domain=0:360, samples=2, samples y=40] ({#4}, {x*(#6)*cos(y)}, {x*(#6)*sin(y)});
  \addplot3[#9, thick, domain=#3:#4, samples=40, samples y=1] ({x}, {0}, {#2});
  \draw[ax3] (axis cs:{min(0,#3)-0.3},0,0) -- (axis cs:#7,0,0);
  \draw[ax3] (axis cs:0,0,{-0.15*(#8)}) -- (axis cs:0,0,#8);
\end{axis}
\end{tikzpicture}}
\newcommand{\solY}[9][0.6]{%
\begin{tikzpicture}[baseline=(current bounding box.center)]
\begin{axis}[mini3d, x={(0.3cm,-0.03cm)}, y={(0.13cm,0.11cm)}, z={(0cm,#1cm)}]
  \addplot3[cap3d, domain=0:1, y domain=0:360, samples=2, samples y=40] ({x*(#5)*cos(y)}, {x*(#5)*sin(y)}, {#3});
  \addplot3[solidsurf, opacity=0.5, point meta=z, domain=#3:#4, y domain=0:360, samples=24, samples y=40] ({(#2)*cos(y)}, {(#2)*sin(y)}, {x});
  \addplot3[cap3d, domain=0:1, y domain=0:360, samples=2, samples y=40] ({x*(#6)*cos(y)}, {x*(#6)*sin(y)}, {#4});
  \addplot3[#9, thick, domain=#3:#4, samples=40, samples y=1] ({#2}, {0}, {x});
  \draw[ax3] (axis cs:0,0,{min(0,#3)-0.15*(#8)}) -- (axis cs:0,0,#8);
  \draw[ax3] (axis cs:-0.3,0,0) -- (axis cs:#7,0,0);
\end{axis}
\end{tikzpicture}}
\newcommand{\minisolo}[1]{\adjustbox{totalheight=1.5cm, max width=4.6cm}{#1}}
\newcommand{\minipair}[2]{\adjustbox{totalheight=1.5cm, max width=4.6cm}{#1\;\raisebox{0.15ex}{\Large$-$}\;#2}}
\newcommand{\minilabel}[1]{\raisebox{0.65cm}{\small #1}\;}


% Small fractions are allowed only in exponents and in the limits of an
% integral (which are superscripts/subscripts); everywhere else use \dfrac.
\newcommand{\dx}{\,dx}
\newcommand{\dy}{\,dy}
\newcommand{\evalat}[3]{\left[\,#1\,\right]_{#2}^{#3}}

% [item 78] right-aligned A-level style marks at the end of a question part
\newcommand{\mk}[1]{{\unskip\nobreak\hfil\penalty50\hskip1.5em\hbox{}\nobreak\hfill\textbf{[#1]}%
  \parfillskip=0pt\finalhyphendemerits=0\par}}
% [item 11] visually emphasised step labels in worked solutions
\newcommand{\step}[1]{\textcolor{teal!55!black}{\textbf{Step #1}}\quad}
% (a), (b), (c) part lists inside questions and solutions
\newenvironment{parts}{\begin{enumerate}[(a)]\setlength{\itemsep}{0.4ex}}{\end{enumerate}}
% SEAB-style solution layout (approved 10.5.8 template): working on the
% left, region sketch on the right (mini 3D pictures go under the sketch)
\newcommand{\solpair}[3][0.56]{\noindent\begin{minipage}[t]{#1\textwidth}\vspace{0pt}#2\end{minipage}\hfill
  \begin{minipage}[t]{\dimexpr0.98\textwidth-#1\textwidth\relax}\vspace{0pt}\centering #3\end{minipage}\par}
% minus-diagram style (drafts/minus_style_template.tex, items 29/31/62)
\pgfplotsset{minus/.style={axis lines=middle, width=5.0cm, height=4.2cm, ticks=none,
  xlabel={$x$}, ylabel={$y$}, every axis x label/.style={at={(ticklabel* cs:1)},anchor=west},
  every axis y label/.style={at={(ticklabel* cs:1)},anchor=south}, clip=false, axis on top,
  axis line style={gray!70}, every axis label/.append style={gray},
  every node/.append style={text=gray!80!black, font=\footnotesize}}}
\tikzset{hlreg/.style={fill=hl!35, draw=hl, thick}}

% house_fonts.tex: HOUSE COPY of definite_integrals/common/prelim_fonts.tex (unchanged)
% =====================================================================
% prelim_fonts.tex: fonts of the real 9758 prelim papers (Word + MathType; pdffonts on
% ACJC 2026 P1/P2, ACJC 2025 P1, JPJC 2026, RVHS 2026, CJC/DHS/TJC/NYJC/HCI/EJC 2025):
%   text Times New Roman; cover/footer Arial; RI page footer Calibri (here Carlito);
%   maths = MathType: Times New Roman letters/digits (italic variables; upright digits,
%   function names, d, e, f), SymbolMT for the integral, sigma, pi and = + < >,
%   scripts 7 pt on 12 pt (about 58 %), tight operator spacing.
% Candidate C (teacher-approved 13 Oct 2026): the real Times New Roman and Arial
% (ttf-mscorefonts-installer) + URW Standard Symbols PS (fonts-urw-base35, a SymbolMT
% clone with the same Symbol encoding) for the integral, sigma, pi, = + < >;
% TeX Gyre Termes Math only for what MathType draws itself (radicals, stretchy brackets,
% fraction bars, arrows). House style: upright d (\mathrm{d}x), upright e and f.
% Loaded once by di_style / de_core (HW sheets) and prelim_style (Tests, solutions).
% =====================================================================
\ifdefined\prelimfontsloaded\endinput\fi
\def\prelimfontsloaded{1}
\edef\pfrestoreat{\catcode`\noexpand\@=\the\catcode`\@\relax}\makeatletter
% MathType script sizes: 7 pt superscripts on 12 pt text (58 %), 5 pt second level
\DeclareMathSizes{12}{12}{7}{5}
\DeclareMathSizes{11}{11}{6.5}{5}
\DeclareMathSizes{10.95}{10.95}{6.5}{5}
\DeclareMathSizes{10}{10}{6}{5}
\DeclareMathSizes{9}{9}{5.5}{5}
\DeclareMathSizes{14.4}{14.4}{8.5}{6}
\setmainfont{Times New Roman}
\setsansfont{Arial}
\newfontfamily\coverfont{Arial}
\newfontfamily\calibrifont{Carlito}
\usepackage{unicode-math}
\setmathfont{TeX Gyre Termes Math}
\setmathfont[range={up/{latin,Latin,num}}]{Times New Roman}
\setmathfont[range={it/{latin,Latin}}]{Times New Roman Italic}
\setmathfont[range={bfup/{latin,Latin,num}}]{Times New Roman Bold}
\setmathfont[range={"002B,"003D,"003C,"003E}]{Standard Symbols PS}
\newfontfamily\pf@sym{Standard Symbols PS}
\newfontfamily\pf@symi[FakeSlant=0.21]{Standard Symbols PS}
% MathType-like spacing (tighter than TeX around = + -)
\AtBeginDocument{\thickmuskip=2.5mu plus 1mu\medmuskip=1.5mu plus 1mu\thinmuskip=3mu
% Symbol integral and sigma (18 pt at 12 pt text, as MathType), slanted Symbol pi
 \renewcommand{\int}{\mathop{\vcenter{\hbox{\pf@sym\fontsize{\fpeval{1.5*\f@size}}{0pt}\selectfont\char"F2}}}\nolimits}%
 \renewcommand{\sum}{\mathop{\vcenter{\hbox{\pf@sym\fontsize{\fpeval{1.5*\f@size}}{0pt}\selectfont\char"E5}}}\limits}%
 \renewcommand{\pi}{\mathord{\text{\pf@symi\char"70}}}}
\pfrestoreat
% [HOUSE 2026-11-09] Times New Roman + TeX Gyre Termes Math
\newif\ifteacher\teachertrue
\ifx\conform\ShowSolutionsKey\else\teacherfalse\fi
\pagestyle{fancy}
\fancyhf{}
% [HOUSE 2026-11-09] chapter header as the Definite Integrals HW sheets (house_format/SPEC.md)
\lhead{\small Definite Integrals: Examples Worksheet\ifteacher\ Solutions\fi}
\rhead{\small Compiled by Maximus Pung Jun}
\cfoot{\small\thepage}
\renewcommand{\headrulewidth}{0.4pt}
\newcounter{wsq}
% \Q{source}: question heading, source cited on the right
% \Q{section tag}{source}: teacher copy starts each question on a new page;
% student copy runs the questions on, separated by a thin rule.
\usepackage{needspace}
\usetikzlibrary{patterns}
\newcommand{\Q}[2]{\ifteacher\ifnum\value{wsq}>0 \clearpage\else\par\bigskip\fi\else\ifnum\value{wsq}>0 \par\bigskip\Needspace{12\baselineskip}\noindent\textcolor{black!25}{\rule{\textwidth}{0.4pt}}\par\medskip\fi\fi
  \stepcounter{wsq}\noindent{\large\bfseries \thewsq}\quad{\small\textcolor{teal!50!black}{\textbf{#1}}}\hfill{\small\textcolor{black!60}{[#2]}}\par\smallskip}% [HOUSE] number only, as \Qhead
% solution box: filled (teacher) or blank, filling the rest of the page (student)
\newtcolorbox{solboxT}{enhanced, width=\dimexpr\linewidth-1.5pt\relax, breakable, colback=white, colframe=black!70, boxrule=0.8pt,
  sharp corners, left=3mm, right=3mm, top=2mm, bottom=2mm,
  title={Solution}, fonttitle=\bfseries\small, coltitle=black, colbacktitle=black!8}
\newtcolorbox{solboxS}{enhanced, width=\dimexpr\linewidth-1.5pt\relax, height fill, colback=white, colframe=black!70, boxrule=0.8pt,
  sharp corners, title={Working}, fonttitle=\bfseries\small, coltitle=black, colbacktitle=black!8}
\newcommand{\Sol}[1]{\ifteacher\par\medskip\begin{solboxT}#1\end{solboxT}\fi}

% ---------------------------------------------------------------------
% [diagram fixes 2026-10-09] REVIEW-COPY SWITCH for the worksheet.
% DI_worksheet_REVIEW.tex defines \wsreviewcopy as 1; the teacher and
% student builds leave it at 0, so \wschg prints its content unchanged.
%   \wschg{note}{diagram}: light-yellow frame behind the diagram + a PDF
%   hover note (Square annotation, written with an xdvipdfmx \special).
% ---------------------------------------------------------------------
\providecommand\wsreviewcopy{0}
\ifnum\wsreviewcopy=1
  \rhead{\textcolor{orange!80!black}{\textbf{REVIEW COPY}} (yellow = not from grail.moe)}
  \lhead{H2 Mathematics 9758 \quad Ch 10}
  \definecolor{wschgyellow}{RGB}{255,246,170}
  \newcommand\wschgann[3]{\special{pdf:ann width #1 height #2 depth 0pt
    << /Type/Annot /Subtype/Square /Contents (#3) /T (Review: examples worksheet)
       /C [1 0.8 0] /BS << /W 0 >> /F 4 >>}}
  \newcommand\wschg[2]{\tcbox[enhanced, on line, colback=wschgyellow!60, colframe=wschgyellow!70!orange,
      boxrule=0.6pt, arc=0pt, boxsep=1pt, left=2pt, right=2pt, top=2pt, bottom=2pt,
      overlay={\path let \p1=(frame.south west), \p2=(frame.north east) in
        node[anchor=south west, inner sep=0, outer sep=0] at (\p1)
        {\wschgann{\the\dimexpr\x2-\x1\relax}{\the\dimexpr\y2-\y1\relax}{#1}};}]{#2}}
  % (worksheet's per-page 'global' note and \Q override removed for the exam practice)
  % [2026-10-10] block version for changed WORKING (displays allowed inside):
  % light-yellow block + hover note covering it
  \newenvironment{wschgblock}[1]{\par\begin{tcolorbox}[enhanced, breakable, colback=wschgyellow!55,
      colframe=wschgyellow!70!orange, boxrule=0.4pt, arc=0pt, boxsep=0pt,
      left=3pt, right=3pt, top=2pt, bottom=2pt, enlarge left by=-3pt, enlarge right by=-3pt,
      before skip=2pt, after skip=2pt, parbox=false,
      overlay={\path let \p1=(frame.south west), \p2=(frame.north east) in
        node[anchor=south west, inner sep=0, outer sep=0] at (\p1)
        {\wschgann{\the\dimexpr\x2-\x1\relax}{\the\dimexpr\y2-\y1\relax}{#1}};}]}%
    {\end{tcolorbox}\par}
\else
  \newcommand\wschg[2]{#2}
  \newenvironment{wschgblock}[1]{}{}
\fi

% ---------------------------------------------------------------------
% [diagram fixes 2026-10-09] Mini 3D solids with BOTH units adjustable, so a
% solid can keep the proportions of its region graph. Same look as \solX /
% \solY in ws_figstyles.tex (which are unchanged).
%  \solXs{x unit cm}{radial unit cm} {R(x)}{a}{b}{R(a)}{R(b)}{x-axis end}{y-axis top}{colour}
%  \solYs{radial unit cm}{y unit cm} {R(y)}{c}{d}{R(c)}{R(d)}{x-axis end}{y-axis top}{colour}
% ---------------------------------------------------------------------
\newcommand{\solXs}[2]{\def\wsxu{#1}\def\wszu{#2}\solXsInner}
\newcommand{\solXsInner}[8]{%
\begin{tikzpicture}[baseline=(current bounding box.center)]
\begin{axis}[mini3d, x={(\wsxu cm,-0.06cm)}, y={({0.15*\wszu cm},{0.1*\wszu cm})}, z={(0cm,\wszu cm)}]
  \addplot3[cap3d, domain=0:1, y domain=0:360, samples=2, samples y=40] ({#2}, {x*(#4)*cos(y)}, {x*(#4)*sin(y)});
  \addplot3[solidsurf, opacity=0.5, point meta=z, domain=#2:#3, y domain=0:360, samples=30, samples y=40] ({x}, {(#1)*cos(y)}, {(#1)*sin(y)});
  \addplot3[cap3d, domain=0:1, y domain=0:360, samples=2, samples y=40] ({#3}, {x*(#5)*cos(y)}, {x*(#5)*sin(y)});
  \addplot3[#8, thick, z buffer=none, domain=#2:#3, samples=40, samples y=1] ({x}, {0}, {#1});
  \draw[ax3] (axis cs:{min(0,#2)-0.15},0,0) -- (axis cs:#6,0,0);
  \draw[ax3] (axis cs:0,0,{-0.15*(#7)}) -- (axis cs:0,0,#7);
\end{axis}
\end{tikzpicture}}
\newcommand{\solYs}[2]{\def\wsxu{#1}\def\wszu{#2}\solYsInner}
\newcommand{\solYsInner}[8]{%
\begin{tikzpicture}[baseline=(current bounding box.center)]
\begin{axis}[mini3d, x={(\wsxu cm,{-0.1*\wsxu cm})}, y={({0.43*\wsxu cm},{0.37*\wsxu cm})}, z={(0cm,\wszu cm)}]
  \addplot3[cap3d, domain=0:1, y domain=0:360, samples=2, samples y=40] ({x*(#4)*cos(y)}, {x*(#4)*sin(y)}, {#2});
  \addplot3[solidsurf, opacity=0.5, point meta=z, domain=#2:#3, y domain=0:360, samples=41, samples y=40] ({(#1)*cos(y)}, {(#1)*sin(y)}, {x});
  \addplot3[cap3d, domain=0:1, y domain=0:360, samples=2, samples y=40] ({x*(#5)*cos(y)}, {x*(#5)*sin(y)}, {#3});
  \addplot3[#8, thick, z buffer=none, domain=#2:#3, samples=80, samples y=1] ({#1}, {0}, {x});
  \draw[ax3] (axis cs:0,0,{min(0,#2)-0.15*(#7)}) -- (axis cs:0,0,#7);
  \draw[ax3] (axis cs:-0.3,0,0) -- (axis cs:#6,0,0);
\end{axis}
\end{tikzpicture}}

% ---------------------------------------------------------------------
% [2026-10-10] Source of the worked solutions.
%   \solsrc{url}: grey line at the top of a solution box citing the school's
%                 own worked solution on grail.moe that the working follows.
%   \pipworked  : grey tag + PDF hover note "Pip-worked, not from grail.moe",
%                 placed on any part whose working is NOT from grail.moe.
% ---------------------------------------------------------------------
\newcommand\solsrc[1]{{\par\noindent\scriptsize\color{black!55}Working follows the school's solution on grail.moe: \url{#1}\par\smallskip}}
\newcommand\pipann[3]{\special{pdf:ann width #1 height #2 depth 0pt
    << /Type/Annot /Subtype/Square /Contents (#3) /T (Solution source)
       /C [0.55 0.55 0.55] /BS << /W 0 >> /F 4 >>}}
\newcommand\pipworkedplain{\begingroup\setbox0\hbox{\colorbox{black!8}{\scriptsize\sffamily Pip-worked}}%
  \edef\pipw{\the\wd0}\edef\piph{\the\ht0}\pipann{\pipw}{\piph}{Pip-worked, not from grail.moe}\box0\endgroup}
\ifnum\wsreviewcopy=1
  \newcommand\pipworked{\wschg{Pip-worked, not from grail.moe}{\scriptsize\sffamily Pip-worked, not from grail.moe}}
\else
  \let\pipworked\pipworkedplain
\fi
